\chapter*{Preface}
\addcontentsline{toc}{chapter}{Preface}
Preface The purpose of this book is to present a substantial collection of the most efficient algorithms for calculating rigid-body dynamics, and to explain them in enough detail that the reader can understand how they work, and how to adapt them (or create new algorithms) to suit the reader’s needs. The collection includes the following well-known algorithms: the recursive Newton-Euler algorithm, the composite-rigid-body algorithm and the articulated-body algorithm. It also includes algorithms for kinematic loops and floating bases. Each algorithm is derived from first principles, and is presented both as a set of equations and as a pseudocode program, the latter being designed for easy translation into any suitable programming language. This book also explains some of the mathematical techniques used to formulate the equations of motion for a rigid-body system. In particular, it shows how to express dynamics using six-dimensional (6D) vectors, and it explains the recursive formulations that are the basis of the most efficient algorithms. Other topics include: how to construct a computer model of a rigid-body system; exploiting sparsity in the inertia matrix; the concept of articulated-body inertia; the sources of rounding error in dynamics calculations; and the dynamics of physical contact and impact between rigid bodies. Rigid-body dynamics has a tendency to become a sea of algebra. However, this is largely the result of using 3D vectors, and it can be remedied by using a 6D vector notation instead. This book uses a notation based on spatial vectors, in which the linear and angular aspects of rigid-body motion are combined into a unified set of quantities and equations. The result is typically a fourto sixfold reduction in the volume of algebra. The benefit is also felt in the computer code: shorter, clearer programs that are easier to read, write and debug, but are still just as efficient as code using standard 3D vectors. This book is intended to be accessible to a wide audience, ranging from senior undergraduates to researchers and professionals. Readers are assumed to have some prior knowledge of rigid-body dynamics, such as might be obtained from an introductory course on dynamics, or from reading the first few chapters of an introductory text. However, no prior knowledge of 6D vectors is required, as this topic is explained from the beginning in Chapter 2. This book does also contain some advanced material, such as might be of interest to dynamics

experts and scholars of 6D vectors. No software is distributed with this book, but readers can obtain source code for most of the algorithms described here from the author’s web site. This text was originally intended to be a second edition of a book entitled Robot Dynamics Algorithms, which was published back in 1987; but it quickly became clear that there was enough new material to justify the writing of a whole new book. Compared with its predecessor, the most notable new materials to be found here are: explicit pseudocode descriptions of the algorithms; a chapter on how to model rigid-body systems; algorithms to exploit branchinduced sparsity; an enlarged treatment of kinematic loops and floating-base systems; planar vectors (the planar equivalent of spatial vectors); numerical errors and model sensitivity; and guidance on how to implement spatial-vector arithmetic.
